% SPDX-License-Identifier: GPL-3.0-or-later
% Copyright (C) 2026 Alessandro Greco
%% EPUB 无障碍替代文本。
\altfig{实验 19 的两条累计曲线：教学模拟包含 2000 名客户端、180 天、每人每天 20 条电路。攻击者分别控制 5\% 的守卫和出口加权带宽。模型中的失守指至少一次使用了守卫和出口均由攻击者控制的电路；到第 180 天，每条电路换守卫的比例为 100\%，持久守卫为 18.4\%。各客户端按模拟的 48--60 天间隔分别轮换守卫，并非同步发生。这不是实际网络去匿名化观测结果。}
% 各客户端的轮换相互独立（模拟间隔 48--60 天），不能标记同步日期。
\begin{tikzpicture}[font=\sffamily\scriptsize,x=0.04cm,y=0.04cm]
  \draw[->] (0,0) -- (190,0) node[below left] {天};
  \draw[->] (0,0) -- (0,108) node[above] {模型失守客户端（\%）};
  \foreach \x in {0,30,60,90,120,150,180} \draw (\x,0) -- (\x,-2) node[below] {\x};
  \foreach \y in {0,25,50,75,100} \draw (0,\y) -- (-2,\y) node[left] {\y};
  \draw[thick,avviso] plot coordinates {(0,0.0) (5,21.1) (10,38.5) (15,51.1) (20,61.5) (25,69.7) (30,76.5) (35,81.7) (40,86.0) (45,89.2) (50,91.7) (55,93.3) (60,94.9) (65,96.2) (70,97.2) (75,97.9) (80,98.2) (85,98.7) (90,98.9) (95,99.1) (100,99.5) (110,99.7) (120,99.9) (165,100.0) (180,100.0)};
  \draw[thick,torviola] plot coordinates {(0,0.0) (1,2.6) (3,4.3) (5,4.4) (48,4.4) (52,5.8) (56,7.1) (60,8.7) (64,9.4) (100,9.5) (108,10.8) (114,12.4) (120,13.2) (130,13.4) (160,14.4) (165,15.6) (170,17.0) (175,18.2) (180,18.4)};
  \node[avviso,anchor=west] at (40,80) {每条电路使用新守卫};
  \node[torviola,anchor=west] at (66,26) {持久守卫};
% 各客户端的轮换相互独立（模拟间隔 48--60 天），不能标记同步日期。
  \node[torgrigio,anchor=west,align=left] at (66,43) {守卫轮换\\并不同步};
\end{tikzpicture}
